\chapter{Kappa-Solutions and Surgery Continuation}

The analytic controls of the preceding chapter produce the ancient models
seen at high curvature.  This chapter classifies those kappa-solutions,
converts their geometry into canonical neighborhoods, and uses a first-failure
argument to restart surgery finitely many times and construct an all-time
controlled flow.

\section{Kappa-Solutions and Asymptotic Models}

\begin{definition}[Kappa-solution]
\label{def:kappa-solution}
\dcref{MT:ch9:9.1,MT:ch9:9.2}
\uses{def:surgery-flow-noncollapsing}
\notready
For $d\in\{2,3\}$, a $d$-dimensional \emph{$\kappa$-solution} is a connected Ricci flow
$(M,g(t))$, $-\infty<t\le0$, such that every time slice is complete and
nonflat, has bounded nonnegative curvature operator, and is
$\kappa$-noncollapsed on every finite scale.  The value of $\kappa>0$ is
fixed for the flow.
\end{definition}

\begin{theorem}[Two-dimensional kappa-solutions are round]
\label{thm:two-dimensional-kappa-solutions-round}
\dcref{MT:ch9:9.4}
\uses{def:kappa-solution,thm:ancient-scalar-curvature-monotonicity}
\notready
Every orientable two-dimensional $\kappa$-solution in the sense of
\ref{def:kappa-solution} is, after a constant
parabolic scaling and time translation, the shrinking round two-sphere.
\end{theorem}

\begin{proof}
The strong maximum principle makes the Gaussian curvature positive: a zero
would make the ancient solution flat, contrary to nonflatness.  A complete
noncompact positively curved surface is diffeomorphic to the plane.  Its
asymptotic circumference ratio is nonincreasing under the flow; if positive,
a blow-down gives a flat cone and Harnack rigidity makes the solution flat,
while if zero, balls at the curvature scale violate noncollapsing.  Hence the
surface is compact.

On a compact surface, write the normalized metric as
$\widehat g(s)=(-t)^{-1}g(t)$ with $s=-\log(-t)$.  Gauss--Bonnet fixes its
average curvature.  The evolution of the variance of the Gaussian curvature,
together with \ref{thm:ancient-scalar-curvature-monotonicity}, is
nonincreasing and has zero backward limit; otherwise a backward blow-down
would be a nonround compact shrinking soliton, contradicting the integrated
soliton equation on a surface.  Thus the curvature is spatially constant on
every slice.  The Ricci-flow equation then gives precisely the shrinking
round sphere.
\end{proof}

\begin{theorem}[Asymptotic models of three-dimensional kappa-solutions]
\label{thm:kappa-solution-asymptotic-models}
\dcref{MT:ch9:9.11,MT:ch9:9.42,MT:ch9:9.44,MT:ch9:9.53,MT:ch9:9.54}
\notready
Assume the endpoint-transfer hypotheses of \ref{thm:reduced-volume-transfer}
hold for the rescaled reduced-exponential families below.  Let $(M,g(t))$ be
a three-dimensional $\kappa$-solution and fix $(p,0)$.
There are $\tau_j\to\infty$ and points $q_j$ at time $-\tau_j$ with reduced
length at most $3/2$ such that the parabolically rescaled pointed flows based
at $(q_j,-1)$ converge to a complete nonflat gradient shrinking soliton with
bounded nonnegative curvature.  The limit is a compact round space form or
one of the three cylindrical models
\[
 S^2\times\mathbb R,\qquad
 \mathbb{RP}^2\times\mathbb R,\qquad
 \mathbb{RP}^2\widetilde\times\mathbb R,
\]
where the last space is the twisted line bundle.  If the original solution
is nonround, a sequence can be chosen with one of these three noncompact
limits.
\end{theorem}

\begin{proof}
\uses{def:kappa-solution,def:reduced-length,def:reduced-volume,lem:reduced-length-support-inequality,lem:minimum-reduced-length-comparison,lem:reduced-density-monotonicity,thm:reduced-volume-transfer,thm:local-ricci-flow-derivative-estimates,thm:hamilton-compactness-generalized-flows,thm:ancient-scalar-curvature-monotonicity,thm:curvature-null-splitting,thm:two-dimensional-kappa-solutions-round}
The minimum comparison \ref{lem:minimum-reduced-length-comparison} gives a
point $q_j$ with $l(q_j,\tau_j)\le3/2$.  The support inequality
\ref{lem:reduced-length-support-inequality} bounds curvature and reduced
length on every fixed rescaled ball about $q_j$.  Noncollapsing and
\ref{thm:hamilton-compactness-generalized-flows} then give an ancient limit
complete at time zero.  Nonnegative curvature passes to the limit, so its
earlier slices are complete as well.  The density formula
\ref{lem:reduced-density-monotonicity} makes its
reduced volume constant in backward time.  Equality in the traced
second-variation inequality forces
the limiting reduced length to be smooth and to satisfy
\[
 \operatorname{Ric}+\operatorname{Hess}l=\frac12g
\]
on the time $-1$ slice.  The strict reduced-volume upper bound excludes a
flat limit.  The local derivative estimates in
\ref{thm:local-ricci-flow-derivative-estimates} and
\ref{thm:ancient-scalar-curvature-monotonicity} give bounded curvature on
each limiting slice.

If the limiting curvature operator has a null direction,
\ref{thm:curvature-null-splitting} splits off a line; the remaining surface
factor is round by \ref{thm:two-dimensional-kappa-solutions-round}.  A free
quotient preserving this splitting is the product, the quotient by the
antipodal map on $S^2$, or the quotient pairing that map with reflection of
the line.  These are precisely the three cylindrical models in the
statement.  If the curvature is positive, tracing and differentiating the soliton equation makes
the potential proper in the noncompact case; its level-set second variation
and Gauss--Bonnet exclude a noncompact positively curved shrinker.  The
compact soliton equation and the maximum principle make the compact case
round.  This gives the listed models.  Finally, a compact asymptotic model
has equality in reduced-volume monotonicity only when the original flow is
the corresponding round shrinker.  A nonround solution therefore admits a
cylindrical asymptotic sequence.
\end{proof}

\begin{proposition}[Universal noncollapsing for nonround kappa-solutions]
\label{prop:universal-nonround-kappa}
\dcref{MT:ch9:9.58}
\notready
Assume the endpoint-transfer hypotheses of \ref{thm:reduced-volume-transfer}
hold uniformly on the rescaled model balls.  There is a universal
$\kappa_*>0$ such that every nonround
three-dimensional $\kappa$-solution is a $\kappa_*$-solution.
\end{proposition}

\begin{proof}
\uses{thm:kappa-solution-asymptotic-models,thm:reduced-volume-transfer}
The asymptotic-model theorem
\ref{thm:kappa-solution-asymptotic-models} gives a backward rescaling which
converges on a fixed parabolic ball to one of the three cylindrical models.
There are only these three models, so a fixed model ball has reduced
volume at least one universal number $V_*>0$.  Smooth convergence gives the
same lower bound, with $V_*/2$, on the approximating rescaled flows.
The quantitative transfer theorem \ref{thm:reduced-volume-transfer} converts
this reduced-volume contribution into a volume ratio
$\kappa_*>0$ on all balls up to the corresponding backward scale.  Undoing
the parabolic rescaling makes that scale tend to infinity.  Shifting an
arbitrary time to zero repeats the argument on every slice, so the same
$\kappa_*$ works at every point and every finite scale.
\end{proof}

\begin{theorem}[Normalized kappa-solution compactness]
\label{thm:kappa-solution-compactness}
\dcref{MT:ch9:9.59,MT:ch9:9.60,MT:ch9:9.62,MT:ch9:9.63,MT:ch9:9.64,MT:ch9:9.65,MT:ch9:9.66,MT:ch9:9.67,MT:ch9:9.68,MT:ch9:9.69,MT:ch9:9.70,MT:ch9:9.71}
\notready
Let $(M_j,g_j(t),p_j)$ be based three-dimensional $\kappa$-solutions with
$R(p_j,0)=1$.  A subsequence converges smoothly on compact spacetime
subsets to a based $\kappa$-solution
\[
 (M_\infty,g_\infty(t),p_\infty)
\]
with $R(p_\infty,0)=1$.
In particular, for every $A<\infty$ and $m\ge0$ there is $C(A,m)<\infty$
such that
\[
 \sup_{B_{g_j(0)}(p_j,A)\times[-A,0]}
 |\nabla^m\operatorname{Rm}_{g_j}|\le C(A,m).
\]
\end{theorem}

\begin{proof}
\uses{def:kappa-solution,thm:hamilton-compactness-generalized-flows,thm:local-ricci-flow-derivative-estimates,thm:local-volume-injectivity-radius-control,thm:curvature-null-splitting,thm:ancient-scalar-curvature-monotonicity,lem:nonflat-terminal-cone-exclusion}
First obtain a uniform scalar-curvature bound on every fixed zero-time ball.
If this failed, choose a nearest point $q_j$ at which the curvature scale is
smaller than its distance scale from $p_j$, and rescale at $q_j$.  The
point-picking choice bounds curvature on a ball whose rescaled radius tends
to infinity.  If the volume ratio of these balls stayed uniformly positive,
$\kappa$-noncollapsing and Hamilton compactness would give a complete
nonflat limit with a zero-curvature point, contradicting
\ref{thm:curvature-null-splitting}.  Otherwise choose the first radius at
which the volume ratio is
half Euclidean.  Volume comparison and the point-picking curvature bound
give a complete flat limit whose unit ball has half the Euclidean volume.
It cannot be Euclidean, while every nontrivial complete flat quotient has
zero asymptotic volume ratio and hence is not noncollapsed on all scales.
This contradiction proves the required based-ball curvature bound.

By \ref{thm:ancient-scalar-curvature-monotonicity}, the same bound holds on
the corresponding backward cylinders.  The local estimates
\ref{thm:local-ricci-flow-derivative-estimates} bound every covariant
derivative on smaller cylinders.
Noncollapsing supplies the base-ball volume lower bound in
\ref{thm:hamilton-compactness-generalized-flows}; diagonalization in the
radius, derivative order, and backward time therefore gives an ancient
smooth limit complete at time zero.  Its nonnegative curvature makes every
earlier slice complete.  It is nonflat because its base scalar curvature is
one, has nonnegative curvature by smooth convergence, and remains
$\kappa$-noncollapsed.

If a limit time slice had unbounded curvature, repeat the point selection
farther and farther from the basepoint.  Blowing down the separating annuli
would produce a nonflat metric cone realized at a positive Ricci-flow time,
contrary to \ref{lem:nonflat-terminal-cone-exclusion}.  Thus every limit
slice has bounded curvature, so the limit satisfies
\ref{def:kappa-solution}.  The derivative bounds obtained above are uniform
before passage to the subsequence and give the last assertion.
\end{proof}

\section{Kappa-Solution Geometry and Stability}

\begin{lemma}[Stability of canonical neighborhoods]
\label{lem:canonical-neighborhood-stability}
\dcref{MT:ch9:9.79,MT:ch9:9.80}
\notready
Fix $\epsilon<1/4$ and $C<\infty$.
\begin{enumerate}
\item A $(C,\epsilon)$-cap, controlled component, or round component
      persists under sufficiently small
      $C^{\lfloor\epsilon^{-1}\rfloor}$ perturbations of its metric and
      under a sufficiently small change of center.
\item An evolving $\epsilon$-neck of normalized backward length strictly
      greater than one persists as a strong $\epsilon$-neck under
      sufficiently small smooth perturbations of the flow.
\item If $C'>C$ and $\epsilon'>\epsilon$, sufficiently close pointed
      metrics carry a $(C',\epsilon')$-canonical neighborhood at the
      corresponding point.
\end{enumerate}
In particular, strong neck centers form an open spacetime set.
\end{lemma}

\begin{proof}
\uses{def:strong-canonical-neighborhood-control,def:canonical-cap,def:epsilon-neck-and-scale}
Curvature ratios, normalized diameters and volumes, core-volume lower
bounds, and curvature-derivative quantities in the definitions all vary
continuously in the stated topology.  A neck chart survives a small metric
perturbation; precomposing it with a small axial translation recenters the
chart.  For a cap, first keep the compact core and its strict estimates,
then truncate and recenter the neck end.  These observations prove the
first assertion and the metric part of the third.

For an evolving neck, the excess normalized time beyond one absorbs the
small change in curvature scale caused by perturbation.  Restricting the
persisting cylinder to normalized length one gives the required strong
neck.  Enlarging $C$ and $\epsilon$ turns all weak boundary inequalities
into strict ones, proving the last metric assertion.  The same argument on
a slightly longer spacetime cylinder proves openness of strong neck centers.
\end{proof}

\begin{proposition}[Curvature trichotomy for kappa-solutions]
\label{prop:kappa-solution-curvature-trichotomy}
\dcref{MT:ch9:9.83}
\notready
Every three-dimensional $\kappa$-solution is exactly one of the following:
it has positive sectional curvature at every time; it is the shrinking round
cylinder $S^2\times\mathbb R$ up to scaling; or it is the free line-bundle
quotient of that cylinder over $\mathbb{RP}^2$.
\end{proposition}

\begin{proof}
\uses{def:kappa-solution,thm:curvature-null-splitting,thm:two-dimensional-kappa-solutions-round}
If the curvature operator has a zero eigenvalue at a positive time, its null
distribution is parallel by \ref{thm:curvature-null-splitting}.  The
universal cover consequently splits isometrically for all earlier times as a
two-dimensional ancient solution times either a line or a circle.  A circle
factor has fixed length while the curved factor expands
backward, so balls at the curvature scale collapse as $t\to-\infty$; this
contradicts $\kappa$-noncollapsing.  The flat factor is therefore a line.
The two-dimensional nonflat factor is the shrinking round sphere by
\ref{thm:two-dimensional-kappa-solutions-round}, and the
only free quotients preserving the product flow are the round cylinder and
the line-bundle quotient over $\mathbb{RP}^2$.  If no zero eigenvalue occurs,
the sectional curvature is positive everywhere.
\end{proof}

\begin{theorem}[Soul structure in nonnegative curvature]
\label{thm:soul-structure-nonnegative-curvature}
\dcref{MT:ch2:2.7}
\notready
Every connected complete noncompact manifold of nonnegative sectional
curvature contains a compact totally geodesic, totally convex soul and is
diffeomorphic to its normal bundle.  If the sectional curvature is positive,
the soul is a point.
\end{theorem}

\begin{proof}
For a ray $\gamma$, the Busemann function
$b_\gamma(x)=\lim_{s\to\infty}(s-d(x,\gamma(s)))$ is concave under
nonnegative curvature by triangle comparison.  Intersect the maximal level
sets of a finite exhaustion by such functions.  Repeating inside the boundary
strata terminates at a compact set $S$ with no boundary.  Concavity makes $S$
totally convex, and the second variation then makes it totally geodesic.
Distance to $S$ has no critical point away from $S$; its gradient flow and
normal exponential map identify $M$ with the normal bundle of $S$.  If
sectional curvature is positive and $S$ contained a nonconstant geodesic, a
flat normal strip obtained from total convexity would contradict positivity.
Thus $S$ is a point.
\end{proof}

\begin{theorem}[Sphere separation and the annulus region about a point soul]
\label{thm:point-soul-sphere-separation-annulus}
\dcref{MT:ch2:2.19}
\uses{thm:soul-structure-nonnegative-curvature}
\notready
Let $M^3$ be complete, noncompact, and positively curved, and let $p$ be its
point soul.  Assume the smooth Schoenflies theorem and the smooth annulus
theorem for embedded two-spheres in the resulting $\mathbb R^3$ identification
of $M\setminus\{p\}$.  Then every smoothly embedded two-sphere in
$M\setminus\{p\}$ has one
relatively compact complementary component with three-ball closure and one
component containing the unique end.  If such a sphere is isotopic in
$M\setminus\{p\}$ to a distance sphere about $p$, then it separates $p$ from
the end.  Any two disjoint embedded two-spheres that separate $p$ from the
end bound a region diffeomorphic to $S^2\times[0,1]$.
\end{theorem}

\begin{proof}
The soul theorem identifies $M\setminus\{p\}$ with the complement of a point
in $\mathbb R^3$.  The assumed smooth Schoenflies theorem gives the ball and
end sides of an embedded sphere, and the assumed annulus theorem identifies
the region between two nested separating spheres with $S^2\times[0,1]$.
Isotopy to a distance sphere preserves which side contains $p$, proving all
three assertions.
\end{proof}

\begin{lemma}[A thin neck separates a point soul from the end]
\label{lem:positive-curvature-neck-separation}
\dcref{MT:ch2:2.19}
\notready
For sufficiently small $\epsilon$, let $M^3$ be complete, noncompact, and
positively curved with point soul $p$.  The central sphere of every
$\epsilon$-neck disjoint from $p$ separates $p$ from the unique end.  Two
disjoint such central spheres bound a region diffeomorphic to
$S^2\times[0,1]$.
\end{lemma}

\begin{proof}
\uses{thm:soul-structure-nonnegative-curvature,lem:metric-sphere-crosses-neck-as-fiber,thm:point-soul-sphere-separation-annulus}
By \ref{thm:soul-structure-nonnegative-curvature}, distance from $p$ has no
critical point away from $p$, so every positive distance sphere is a
two-sphere separating $p$ from the end, and identifies $M$ with
$\mathbb R^3$.

Now take $z=p$ in
\ref{lem:metric-sphere-crosses-neck-as-fiber} and choose any point in the
middle sixth of the neck.  The resulting distance sphere is a graph over,
and hence isotopic inside the neck to, the central sphere.  The sphere
separation theorem \ref{thm:point-soul-sphere-separation-annulus} therefore
puts the soul and the end on opposite sides of the central sphere.  Applying
its annulus conclusion to two disjoint central spheres gives the claimed
$S^2\times[0,1]$ region.
\end{proof}

\begin{proposition}[Positive curvature excludes arbitrarily small necks]
\label{prop:positive-curvature-neck-scale-lower-bound}
\dcref{MT:ch2:2.20}
\notready
For every sufficiently small $\epsilon>0$, a complete positively curved
Riemannian three-manifold contains no sequence of $\epsilon$-necks whose
scales tend to zero.
\end{proposition}

\begin{proof}
\uses{def:epsilon-neck-and-scale,thm:soul-structure-nonnegative-curvature,thm:busemann-function-superharmonicity,lem:epsilon-neck-geodesic-axis-alignment,lem:positive-curvature-neck-separation}
The assertion is immediate on a compact manifold.  Otherwise
\ref{thm:soul-structure-nonnegative-curvature} gives a point soul $p$ and
identifies the manifold with $\mathbb R^3$.  By
\ref{lem:positive-curvature-neck-separation}, disjoint necks away from the
soul can be ordered toward the end.

Choose two such ordered necks $N_1,N_2$, with centers $x_i$ and scalar
curvatures $R_i$.  Let $B$ be the Busemann function of a ray from $p$, and
choose a nonnegative cutoff $\psi$ which is one between the central regions
and changes only along the middle thirds of the two necks.  The alignment
estimate \ref{lem:epsilon-neck-geodesic-axis-alignment} compares $\nabla B$
with the signed axial direction on both transition regions.  Weak
superharmonicity from \ref{thm:busemann-function-superharmonicity} and
integration by parts give
\[
 0\le \int_M\langle\nabla B,\nabla\psi\rangle\,d\mu
   =\bigl(\alpha_2R_2^{-1}-\alpha_1R_1^{-1}\bigr)
      \operatorname{Vol}(S^2,h_0),
\]
where the cylindrical neck estimates give
$|\alpha_i-1|\le c\epsilon$.  After decreasing $\epsilon$ so that
$c\epsilon<1/3$, this implies $R_2\le2R_1$.  Fixing the first neck therefore
bounds below the scale $R^{-1/2}$ of every later neck.  On the compact region
on the soul side, scalar curvature is bounded, so neck scales have a positive
lower bound there as well.  Hence no sequence of neck scales can tend to
zero.
\end{proof}

\begin{theorem}[Splitting at infinity for kappa-solutions]
\label{thm:kappa-solution-splitting-at-infinity}
\dcref{MT:ch9:9.39}
\notready
Let $(M,g(t))$, $-\infty<t\le0$, be a noncompact three-dimensional
$\kappa$-solution.  If $p_j\to\infty$ and
\[
 R(p_j,-1)d_{g(-1)}(p,p_j)^2\longrightarrow\infty,
\]
then there are points $q_j\to\infty$ and scales $Q_j=R(q_j,-1)$ for
which the based parabolic rescalings converge on compact subsets to the
product of a shrinking round two-sphere with a flat line, possibly modulo
the free line-bundle involution.
\end{theorem}

\begin{proof}
\uses{def:kappa-solution,thm:ancient-scalar-curvature-monotonicity,thm:hamilton-compactness-generalized-flows,thm:curvature-null-splitting,thm:two-dimensional-kappa-solutions-round}
Point-pick $q_j$ in the ball of half the distance from $p_j$ to maximize
scalar curvature times the square of the remaining distance.  Then
$Q_jd(p,q_j)^2\to\infty$ and $R\le4Q_j$ on balls whose rescaled radii tend
to infinity.  By \ref{thm:ancient-scalar-curvature-monotonicity} this bound
holds on the corresponding backward cylinders.  Noncollapsing and
\ref{thm:hamilton-compactness-generalized-flows} give a nonflat ancient limit
complete at time zero; nonnegative curvature makes all its earlier slices
complete.

The rescaled minimizing geodesics from $p$ to $q_j$ converge on each compact
subinterval to a line through the basepoint.  The splitting clause of
\ref{thm:curvature-null-splitting} therefore writes the limit as a
two-dimensional ancient flow times a flat line.  The surface factor is
nonflat and noncollapsed, hence round by
\ref{thm:two-dimensional-kappa-solutions-round}.  The only free quotient
compatible with the line factor is the stated line-bundle involution.
\end{proof}

\begin{theorem}[Structure of noncompact kappa-solutions]
\label{thm:noncompact-kappa-solution-structure}
\dcref{MT:ch9:9.84,MT:ch9:9.85,MT:ch9:9.87,MT:ch9:9.88}
\notready
For every sufficiently small $\epsilon>0$ there is
$C_0=C_0(\epsilon)<\infty$ such that the time-zero slice of every
noncompact three-dimensional $\kappa$-solution containing no embedded
two-sided $\mathbb{RP}^2$ is either a strong $\epsilon$-tube or a
$(C_0,\epsilon)$-cap followed by a strong $\epsilon$-tube.  The cap core has
uniform scale-normalized curvature, derivative, diameter, and volume bounds.
\end{theorem}

\begin{proof}
\uses{def:kappa-solution,prop:universal-nonround-kappa,prop:kappa-solution-curvature-trichotomy,thm:soul-structure-nonnegative-curvature,thm:kappa-solution-splitting-at-infinity,thm:kappa-solution-compactness,lem:canonical-neighborhood-stability,def:canonical-neck-cap-regions,thm:canonical-neck-cap-topology}
Consider first the positive-curvature case and normalize at a soul $p$ by
$R(p,0)=1$, using \ref{thm:soul-structure-nonnegative-curvature}.  Suppose
points $q_j$ escaping every fixed ball were not strong
$\epsilon$-neck centers.  If
$d(p,q_j)^2R(q_j,0)\to\infty$,
\ref{thm:kappa-solution-splitting-at-infinity} gives the shrinking round
cylinder, and
\ref{lem:canonical-neighborhood-stability} makes $q_j$ a neck center, a
contradiction.  If those products remain bounded, then
$R(q_j,0)\to0$ while the unit-curvature soul stays at bounded distance in
the $R(q_j,0)$-rescaling.  The based curvature bound in
\ref{thm:kappa-solution-compactness} rules this out.

The universal constant in \ref{prop:universal-nonround-kappa} makes the
radius beyond which all points are neck centers uniform.  Otherwise
normalize a sequence of souls and take the first non-neck point beyond an
ever longer neck annulus.  Compactness at the souls transfers a fixed neck
annulus to the sequence.  At the outer first-failure point the preceding
scale-separation argument again gives a cylindrical limit, contradicting
stability.  Compactness on the remaining soul ball supplies uniform
curvature ratios, derivative bounds, volume bounds, and diameter bounds.
The neck spheres separate the unique end, and
\ref{thm:canonical-neck-cap-topology} identifies the bounded remainder as a
cap core.

In the zero-curvature cases of
\ref{prop:kappa-solution-curvature-trichotomy}, the round cylinder is already
a strong tube.  Its nontrivial orientable quotient has one controlled cap
core and a cylindrical end.  The excluded two-sided
$\mathbb{RP}^2$ removes the remaining exceptional product.  Increasing
$C_0$ to dominate the uniform core estimates proves the theorem.
\end{proof}

\begin{theorem}[Structure of compact kappa-solutions]
\label{thm:compact-kappa-solution-structure}
\dcref{MT:ch9:9.89,MT:ch9:9.90,MT:ch9:9.91}
\notready
For every sufficiently small $\epsilon>0$ there is
$C_1=C_1(\epsilon)<\infty$ such that a compact three-dimensional
$\kappa$-solution is round, is a positively curved $C_1$-component, or is
the union of two disjoint $(C_1,\epsilon)$-cap cores joined by a strong
$\epsilon$-tube.  In the nonround cases its time-zero slice is
diffeomorphic to $S^3$ or $\mathbb{RP}^3$.
\end{theorem}

\begin{proof}
\uses{def:kappa-solution,prop:universal-nonround-kappa,prop:kappa-solution-curvature-trichotomy,thm:kappa-solution-compactness,thm:noncompact-kappa-solution-structure,lem:canonical-neighborhood-stability,thm:canonical-neck-cap-topology}
The splitting alternatives in
\ref{prop:kappa-solution-curvature-trichotomy} are noncompact, so a compact
nonround solution has positive sectional curvature and finite fundamental
group.  Normalize its maximum scalar curvature to one.  If a sequence has
unbounded normalized diameter and contains a point with neither a cap nor a
neck neighborhood, compactness based at that point gives a noncompact
$\kappa$-solution.  The point lies in the cap/tube decomposition of
\ref{thm:noncompact-kappa-solution-structure}, and stability transfers that
neighborhood back, a contradiction.  Hence every large-diameter solution is
a double-capped strong tube.  The topology classification in
\ref{thm:canonical-neck-cap-topology}, together with finite fundamental
group, leaves $S^3$ and $\mathbb{RP}^3$.

For bounded normalized diameter,
\ref{prop:universal-nonround-kappa} puts every nonround member in one fixed
noncollapsed family, which is compact by
\ref{thm:kappa-solution-compactness}.  A nonround limit retains positive
curvature, so sectional-curvature ratios, diameter, volume, and curvature
derivative ratios have positive finite extrema over this family.  Taking
$C_1$ larger than those extrema makes the whole slice a controlled
component.  A round limit gives the round alternative.  These cases exhaust
the compact family and prove the assertion.
\end{proof}

\begin{theorem}[Uniform structure of kappa-solutions]
\label{thm:kappa-solution-structure-summary}
\dcref{MT:ch9:9.93}
\notready
For every sufficiently small $\epsilon>0$ there is
$C=C(\epsilon)<\infty$ such that every three-dimensional
$\kappa$-solution is one of the following: a round spherical space form; a
compact positively curved $C$-component; a compact double-capped strong
$\epsilon$-tube; a noncompact $(C,\epsilon)$-capped strong
$\epsilon$-tube; the round strong cylinder; or
$\mathbb{RP}^2\times\mathbb R$.  In every nonexceptional case each point is
in a cap core, a controlled or round component, or is the center of a strong
$\epsilon$-neck.  Uniformly on the flow,
\[
 |\nabla R|\le CR^{3/2},\qquad |\partial_tR|\le CR^2.
\]
\end{theorem}

\begin{proof}
\uses{prop:universal-nonround-kappa,prop:kappa-solution-curvature-trichotomy,thm:noncompact-kappa-solution-structure,thm:compact-kappa-solution-structure,thm:kappa-solution-compactness}
Apply \ref{prop:kappa-solution-curvature-trichotomy}.  Its two split cases
give the cylinder and the exceptional product.  The positive-curvature
noncompact case is exactly
\ref{thm:noncompact-kappa-solution-structure}, and the compact case is
\ref{thm:compact-kappa-solution-structure}.  Choose $C$ larger than the
constants in both results.  By
\ref{prop:universal-nonround-kappa}, the local derivative estimates applied
after normalizing any nonround point to scalar curvature one are uniform on
the compact
normalized family of \ref{thm:kappa-solution-compactness}; scaling back gives
the two displayed inequalities.  This establishes every alternative and
the uniform constants simultaneously.
\end{proof}

\section{Canonical Neighborhood Models}

\begin{theorem}[Canonical neighborhoods in kappa-solutions]
\label{thm:kappa-solution-canonical-neighborhoods}
\dcref{MT:ch9:9.94,MT:ch9:9.95}
\notready
For every sufficiently small $\epsilon>0$ there is
$C=C(\epsilon)<\infty$ such that every point of every three-dimensional
$\kappa$-solution has a strong $(C,\epsilon)$-canonical neighborhood,
unless the solution is the round product $\mathbb{RP}^2\times\mathbb R$.
Moreover, if generalized flows with no embedded two-sided
$\mathbb{RP}^2$ converge smoothly to a $\kappa$-solution, then their
basepoints have strong $(C,\epsilon)$-canonical neighborhoods for all
sufficiently large indices.
\end{theorem}

\begin{proof}
\uses{def:strong-canonical-neighborhood-control,thm:kappa-solution-structure-summary,lem:canonical-neighborhood-stability}
Every nonexceptional alternative in
\ref{thm:kappa-solution-structure-summary} is one of the four alternatives
in \ref{def:strong-canonical-neighborhood-control}; taking the same uniform
$C(\epsilon)$ proves existence.  The only remaining alternative is the round
$\mathbb{RP}^2\times\mathbb R$ product.

For the convergence assertion, the limiting basepoint has a canonical
neighborhood by the first part.  A cap, controlled component, or round
component transfers directly by the first part of
\ref{lem:canonical-neighborhood-stability}.  A strong neck extends to a
slightly longer normalized backward interval, so the second part of that
lemma transfers it as a strong neck.  The limit cannot be the exceptional
product: an embedded central $\mathbb{RP}^2$ has trivial normal bundle and
persists under smooth convergence, contrary to the hypothesis on the
approximating slices.  Thus the basepoints eventually have the asserted
canonical neighborhoods.
\end{proof}

\begin{theorem}[Canonical neighborhoods in the standard solution]
\label{thm:standard-solution-canonical-neighborhoods}
\dcref{MT:ch12:12.28,MT:ch12:12.32,MT:ch12:12.36}
\notready
For every sufficiently small $\epsilon>0$ there is
$C_{\mathrm{std}}(\epsilon)<\infty$ such that every point $(q,t)$ of the
standard solution has one of the following neighborhoods:
\begin{enumerate}
\item a $(C_{\mathrm{std}},\epsilon)$-cap whose core contains $q$;
\item an evolving $\epsilon$-neck reaching time zero, with its initial slice
      disjoint from the cap core;
\item an evolving $\epsilon$-neck defined for normalized backward time
      $1+\epsilon$.
\end{enumerate}
The alternatives persist for flows converging smoothly to the standard
solution on the corresponding parabolic neighborhoods.
\end{theorem}

\begin{proof}
\uses{def:strong-canonical-neighborhood-control,def:standard-surgery-cap-model,thm:hamilton-compactness-generalized-flows,thm:kappa-solution-canonical-neighborhoods,lem:canonical-neighborhood-stability}
The standard solution is rotationally symmetric, positively curved, and
asymptotic to the evolving round cylinder.  Its scalar curvature tends
uniformly to infinity as $t\uparrow1$.  Normalize any sequence of points with
unbounded curvature.  Noncollapsing and the derivative bounds give a smooth
limit by \ref{thm:hamilton-compactness-generalized-flows}; the limit is a
noncompact $\kappa$-solution, hence cylindrical away from a possible cap.
Thus all sufficiently high-curvature points have either a cap neighborhood
or an evolving neck which extends backward for normalized time
$1+\epsilon$.

Choose $\theta<1$ so that this covers $t\ge\theta$.  On
$[0,\theta]$, asymptotic cylindricality is uniform outside one compact
spacetime set, giving necks which reach time zero.  On the remaining compact
set, positivity of curvature and smoothness give uniform curvature ratios,
diameter, volume, and derivative bounds at the curvature scale.  Enlarging
the fixed core of \ref{def:standard-surgery-cap-model} and then
$C_{\mathrm{std}}$
makes this set a controlled cap.  Smooth pointed convergence preserves each
strict estimate after a harmless enlargement of the constants, proving the
last assertion.
\end{proof}

\section{First-Failure Extension}

\begin{lemma}[First canonical-neighborhood failure]
\label{lem:canonical-neighborhood-first-failure}
\dcref{MT:ch17:17.2}
\notready
Fix a finite induction stage and a sequence $r_a\downarrow0$.  Suppose that
for every $a$ there are controlled surgery flows with surgery controls
tending to zero for which the strong $(C,\epsilon)$-canonical-neighborhood
assumption with parameter $r_a$ fails before the stage endpoint.  For each
$a$, after passing sufficiently far down that sequence, there is a first
failure time $t_{a,b}$ and a point $x_{a,b}$ at that time such that
$R(x_{a,b})\ge r_a^{-2}$, the assumption holds before $t_{a,b}$, and no
strong $(C,\epsilon)$-canonical neighborhood exists at $x_{a,b}$.
\end{lemma}

\begin{proof}
\uses{def:controlled-ricci-flow-with-surgery,lem:canonical-neighborhood-stability,thm:standard-solution-canonical-neighborhoods,thm:local-neck-surgery,lem:local-surgery-standard-cap-convergence}
For fixed $a,b$, take failure points whose times decrease to the infimum of
all failure times.  The preceding time slice is compact, so after flowing
the points to a common nearby time and taking a subsequence they converge to
a point $x_{a,b}$.  Curvature continuity gives
$R(x_{a,b})\ge r_a^{-2}$.  Caps, controlled components, and round components
are open under smooth time convergence by
\ref{lem:canonical-neighborhood-stability}; if the limiting point had one of
these neighborhoods, the approximating failure points would have the same
alternative.

Suppose instead that the limiting point is the center of a strong neck.
If its backward cylinder remains in the continuing region, the spacetime
part of \ref{lem:canonical-neighborhood-stability} again transfers the neck
to the approximating points.  Otherwise that cylinder meets a newly inserted
cap.  Rescale by the surgery scale.  The local construction and
\ref{lem:local-surgery-standard-cap-convergence} identify a larger and
larger parabolic ball with the standard solution.  The alternatives in
\ref{thm:standard-solution-canonical-neighborhoods} give either a cap around
the failure point or an evolving neck overlapping the original surgery
neck; the overlap charts from \ref{thm:local-neck-surgery} glue the latter to
a strong neck.  Both alternatives contradict failure.  Hence the infimum is
attained at a genuine failure point with the stated preceding control.
\end{proof}

\begin{lemma}[A bounded-curvature cap encounter is standard]
\label{lem:standard-cap-encounter-control}
\dcref{MT:ch17:17.4,MT:ch17:17.5,MT:ch17:17.8}
\notready
Consider rescaled first-failure flows and a backward flow line, starting in
a fixed-radius zero-time ball, which reaches a surgery cap after rescaled
time $s_a\le s'<\infty$.  If scalar curvature along the line is bounded by
$D'<\infty$ and $\bar h_a$ is the rescaled surgery scale, then for all large
$a$ there are constants $c=c(D',s')>0$, $A<\infty$, and $\theta<1$ such that
\[
 \bar h_a^2\ge c,
 \qquad s_a<\theta\bar h_a^2,
\]
and the point where the line reaches the surgery slice lies in the image of
the standard-solution ball $B(p_0,A)$ used to model the cap.
\end{lemma}

\begin{proof}
\uses{def:controlled-surgery-parameters,thm:local-neck-surgery,lem:local-surgery-standard-cap-convergence,thm:standard-solution-canonical-neighborhoods}
At the cap boundary, the surgery construction compares scalar curvature to
$\bar h_a^{-2}$ within a fixed factor $D$.  The assumed bound $D'$ along the
flow line therefore gives
\[
 \bar h_a^2\ge(2D'D)^{-1}.
\]
To reach the inserted cap, the flow line first crosses the retained collar
of the cutting neck.  In the neck coordinates of
\ref{thm:local-neck-surgery}, that collar has fixed normalized axial length,
and the local surgery construction places its inner boundary a fixed
multiple of $\bar h_a$ from the cap tip.  Thus the entry point is at distance
at most $C\bar h_a$ from the tip.  After division by $\bar h_a$, it lies in
one fixed standard ball $B(p_0,A)$.

It remains to bound elapsed time away from the standard singular time.  If
$s_a\le\bar h_a^2/2$, any $\theta\in(1/2,1)$ works.  Otherwise the curvature
bound along the line gives a fixed upper bound after cap-scale rescaling.
Choose $\theta<1$ so close to one that the scalar curvature of the standard
solution on the corresponding cap ball is larger than that bound.  Smooth
cap convergence then contradicts $s_a\ge\theta\bar h_a^2$.  This proves all
three asserted controls.
\end{proof}

\begin{lemma}[A cap-reaching flow line forces a canonical neighborhood]
\label{lem:bounded-flow-line-cap-exclusion}
\dcref{MT:ch17:17.7}
\notready
In a rescaled first-failure sequence, suppose a point in a fixed-radius
zero-time ball has a backward flow line of uniformly bounded scalar
curvature which reaches a surgery cap in uniformly bounded time.  Then the
base first-failure points have strong $(C,\epsilon)$-canonical neighborhoods
for all sufficiently large indices.
\end{lemma}

\begin{proof}
\uses{lem:standard-cap-encounter-control,lem:local-surgery-standard-cap-convergence,thm:standard-solution-canonical-neighborhoods,lem:canonical-neighborhood-stability,thm:local-neck-surgery}
By \ref{lem:standard-cap-encounter-control}, the cap scale is bounded below,
the elapsed cap-scale time is at most some $\theta<1$, and the relevant
backward point lies in one fixed standard-solution ball.  Enlarge that ball
so that forward cap evolution contains the basepoint at time zero.  Smooth
convergence from \ref{lem:local-surgery-standard-cap-convergence} then
identifies this entire parabolic region with the corresponding standard
solution region.

Apply \ref{thm:standard-solution-canonical-neighborhoods} at the limiting
basepoint.  A cap alternative transfers directly by
\ref{lem:canonical-neighborhood-stability}.  In an evolving-neck
alternative, the neck reaches the continuing side of the original surgery
neck; their overlap charts agree after a small axial adjustment, and
\ref{thm:local-neck-surgery} glues them into a strong
$\epsilon$-neck.  Thus either alternative gives a forbidden canonical
neighborhood at the first-failure point.
\end{proof}

\begin{lemma}[High-curvature control in a rescaled failure slice]
\label{lem:rescaled-high-curvature-canonical-control}
\dcref{MT:ch17:17.6}
\notready
After rescaling first-failure points to scalar curvature one and taking the
surgery controls to zero, every point of the zero-time slice with scalar
curvature greater than one has a strong
$(2C,2\epsilon)$-canonical neighborhood for all sufficiently large indices.
\end{lemma}

\begin{proof}
\uses{lem:canonical-neighborhood-first-failure,lem:canonical-neighborhood-stability,lem:local-surgery-standard-cap-convergence,thm:standard-solution-canonical-neighborhoods}
A nonexposed point is approached along its backward flow line by points at
times strictly before the first failure.  At those points the original
$(C,\epsilon)$-canonical-neighborhood assumption applies because their
scalar curvature is above the rescaled threshold.  Smooth convergence to
time zero and Part~(3) of
\ref{lem:canonical-neighborhood-stability} transfer a
$(2C,2\epsilon)$ neighborhood.

An exposed point lies in a newly inserted cap.  As surgery control tends to
zero, \ref{lem:local-surgery-standard-cap-convergence} identifies its
bounded cap-scale neighborhood with the standard solution, and
\ref{thm:standard-solution-canonical-neighborhoods} supplies a cap or strong
neck alternative.  Stability transfers it with doubled constants.  These
cases exhaust the zero-time slice.
\end{proof}

\begin{lemma}[Positive components at first failure]
\label{lem:positive-component-first-failure-control}
\dcref{MT:ch16:16.1,MT:ch17:17.1}
\uses{def:strong-canonical-neighborhood-control,def:controlled-ricci-flow-with-surgery,thm:conditional-noncollapsing-extension,thm:bounded-curvature-at-bounded-distance,thm:pointed-controlled-flow-compactness,thm:kappa-solution-canonical-neighborhoods,lem:bounded-flow-line-cap-exclusion,lem:rescaled-high-curvature-canonical-control,lem:canonical-neighborhood-stability}
\notready
In a rescaled first-failure sequence at a fixed positive induction stage,
with unit scalar curvature at the basepoints, suppose the component
containing each basepoint has positive sectional curvature.  Assume in
addition that the sequence is ordinarily noncollapsed on every fixed based
ball, pointed compactness produces a complete nonflat limit, and every
positive-curvature limit satisfies the stated kappa-solution canonical
neighborhood theorem.  Then, after passing to a subsequence, the basepoints
eventually have strong
$(C,\epsilon)$-canonical neighborhoods.
\end{lemma}

\begin{lemma}[First-failure noncollapsing bridge]
\label{lem:first-failure-noncollapsing-bridge}
\dcref{MT:ch16:16.1,MT:ch17:17.1}
\notready
Consider a rescaled first-failure sequence with unit scalar curvature at the
basepoints and with surgery controls tending to zero.  After passing to a
subsequence, one of the following holds:
\begin{enumerate}
\item the basepoints eventually have strong
      $(C,\epsilon)$-canonical neighborhoods;
\item there is $\kappa_*>0$ such that, on every fixed zero-time based ball,
      every sufficiently advanced flow satisfies the ordinary
      $\kappa_*$-noncollapsing volume estimate on each fixed scale on which
      the required backward parabolic neighborhood exists.
\end{enumerate}
\end{lemma}

\begin{proof}
\uses{def:surgery-flow-noncollapsing,thm:conditional-noncollapsing-extension,lem:positive-component-first-failure-control}
Conditional noncollapsing,
Theorem~\ref{thm:conditional-noncollapsing-extension}, supplies one constant
on every ball whose center is not in a positively curved component.  Thus
only the positive-component alternative needs to be removed.  If it occurs
at a point of a fixed based ball, connectedness puts the unit-curvature
basepoint in the same compact positive component.  The positive-component
lemma \ref{lem:positive-component-first-failure-control} then gives the
first alternative.

Consequently a subsequence whose basepoints remain genuine failures contains
no positive-component exception on any fixed based ball.  Diagonalizing in
the radius leaves exactly the ordinary volume estimate with the stage
constant, which is the second alternative.
\end{proof}

\begin{lemma}[Short backward curvature control]
\label{lem:short-backward-curvature-control}
\dcref{MT:ch17:17.9}
\notready
For every $A<\infty$ there are $D(A)<\infty$ and $\sigma(A)>0$ such that,
in every sufficiently advanced rescaled first-failure flow, scalar curvature
is at most $D(A)$ along each backward flow line starting in
$B(x_{a,b},0,A)$ for as long as that line exists in
$[-\sigma(A),0]$.
\end{lemma}

\begin{proof}
\uses{lem:first-failure-noncollapsing-bridge,thm:bounded-curvature-at-bounded-distance,lem:rescaled-high-curvature-canonical-control}
The canonical alternative in
\ref{lem:first-failure-noncollapsing-bridge} contradicts the choice of
first-failure basepoints.  Its other alternative therefore gives one
ordinary noncollapsing constant on every fixed based ball.  By
\ref{lem:rescaled-high-curvature-canonical-control}, every zero-time point
of scalar curvature above one has doubled canonical control.  The
bounded-curvature-at-bounded-distance argument therefore applies from the
unit-curvature basepoint and gives a bound $D_0(A)$ on the entire zero-time
ball.

On a canonical neck, cap, or controlled component, the scale-invariant
estimate $|\partial_tR|\le cR^2$ holds with a uniform $c$; it also holds on
the compact low-curvature remainder by the Ricci-flow local derivative
estimate.  Along a backward flow line starting with
$R\le D_0(A)$, integration of this differential inequality bounds
$R$ by $2D_0(A)$ for a time depending only on $A$.  Pinching converts this
to a full curvature bound.  Taking $D(A)=2D_0(A)$ and the corresponding
$\sigma(A)$ proves the statement.
\end{proof}

\begin{lemma}[Backward parabolic-neighborhood dichotomy]
\label{lem:backward-parabolic-neighborhood-dichotomy}
\dcref{MT:ch17:17.10}
\notready
After passing to a subsequence of rescaled first-failure flows, exactly one
of the following occurs:
\begin{enumerate}
\item for every $A<\infty$ there are $D(A),s(A)>0$ such that the compatible
      parabolic neighborhood
      $P(x_{a,b},0,A,-s(A))$ exists and has
      $|\operatorname{Rm}|\le D(A)$;
\item eventually every basepoint has a strong
      $(C,\epsilon)$-canonical neighborhood.
\end{enumerate}
\end{lemma}

\begin{proof}
\uses{lem:short-backward-curvature-control,lem:bounded-flow-line-cap-exclusion}
Fix an integer radius $A$.  The bound in
\ref{lem:short-backward-curvature-control} controls every backward flow line
from the zero-time $A$-ball for a uniform short time.  If all such lines
survive, their union is the required compatible parabolic neighborhood.  If
one does not survive, it reaches a surgery boundary with uniformly bounded
curvature and in uniformly bounded time, so
\ref{lem:bounded-flow-line-cap-exclusion} gives a canonical neighborhood at
the basepoint.  Diagonalize over integer $A$.  Either the second outcome
occurs along a subsequence, or the first outcome holds for every radius.
\end{proof}

\begin{lemma}[Bounded curvature of a first-failure limit slice]
\label{lem:first-failure-limit-slice-bounded-curvature}
\dcref{MT:ch11:11.7}
\notready
Let a time slice of a pointed first-failure limit be complete, nonflat, and
nonnegatively curved.  Suppose it is ordinarily $\kappa$-noncollapsed on all
finite scales and is the smooth limit of slices on which points above a fixed
scalar-curvature threshold have uniform canonical neighborhoods.  Then its
curvature is bounded on the whole slice.
\end{lemma}

\begin{proof}
\uses{def:surgery-flow-noncollapsing,thm:curvature-null-splitting,lem:rescaled-high-curvature-canonical-control,prop:positive-curvature-neck-scale-lower-bound}
If the curvature operator has a null eigenvalue, the strong maximum principle
and the local clause of \ref{thm:curvature-null-splitting}, applied on every
bounded parabolic ball before passing to the exhaustion, give a compatible
local product by a positively curved surface and a one-manifold.  These
products globalize on the universal cover, by the complete-slice clause of
that theorem, to $\Sigma\times\mathbb R$.  If scalar curvature never exceeds the canonical-
neighborhood threshold, it is already bounded.  Otherwise choose a point
above that threshold and lift it to the product cover.  Its lifted canonical
neighborhood is a neck or contains a
neck; whole compact-component alternatives cannot cover
$\Sigma\times\mathbb R$.  On this sufficiently thin neck the unique tangent
Ricci-null line is uniformly close to the axial line: its orthogonal planes
carry the only positive curvature eigenvalue.  The central neck sphere is
therefore transverse to the parallel line foliation and projects as a local
covering onto the surface factor.  Its image is both open and compact, hence
all of connected $\Sigma$.  Thus $\Sigma$ is compact, its continuous
curvature is bounded, and so is the curvature of the product slice and its
quotient.
If the slice is compact, boundedness follows directly by continuity.

It remains to consider a noncompact slice with positive sectional curvature.
Suppose its curvature were unbounded.  Choose points escaping compact sets
with scalar curvature tending to infinity.  The smooth-limit form of
\ref{lem:rescaled-high-curvature-canonical-control} gives a doubled canonical
neighborhood at every sufficiently high-curvature point.  A round-component
or $C$-component alternative is impossible because the component is
noncompact.  A cap alternative contains, at a scale comparable to its core
scale, the neck forming its noncompact end.  Thus every selected canonical
neighborhood contains a $2\epsilon$-neck whose scale is comparable to
$R^{-1/2}$ and hence tends to zero.  This contradicts
\ref{prop:positive-curvature-neck-scale-lower-bound}.  Thus the curvature of
the complete slice is globally bounded.
\end{proof}

\begin{lemma}[Uniform backward lifetime]
\label{lem:uniform-backward-lifetime}
\dcref{MT:ch17:17.11}
\uses{lem:backward-parabolic-neighborhood-dichotomy,lem:first-failure-limit-slice-bounded-curvature,lem:short-backward-curvature-control,lem:rescaled-high-curvature-canonical-control,thm:hamilton-compactness-generalized-flows,lem:bounded-flow-line-cap-exclusion}
\notready
Suppose the first alternative of
\ref{lem:backward-parabolic-neighborhood-dichotomy} holds and the basepoints
remain first-failure points.  The backward lifetimes $s(A)$ may then be
chosen bounded below by one number $s_*>0$ independent of $A$.  Extending
maximally, either the compatible neighborhoods exist for every backward
time, or a flow line starting in one fixed-radius zero-time based ball
reaches a surgery cap at the finite maximal backward time with curvature
bounded in the original first-failure normalization.
\end{lemma}

\begin{theorem}[Extension of strong canonical neighborhoods]
\label{thm:strong-canonical-neighborhood-extension}
\dcref{MT:ch17:17.1}
\notready
Assume the positive-component first-failure, endpoint-transfer, pointed-limit,
and kappa-solution canonical-neighborhood interfaces used in the proof hold
with constants uniform in the stage parameters.
Under the finite stage-$i$ parameter hypotheses of
\ref{thm:conditional-noncollapsing-extension}, there are
$0<r_{i+1}\le r_i$ and
$0<\delta_{i+1}\le\delta(r_{i+1})$ such that the following holds.  If a
controlled Ricci flow with surgery on $[0,T)$,
$T\in(T_i,T_{i+1}]$, has $\bar\delta\le\delta_{i+1}$ on
$[T_{i-1},T]$ and satisfies the stage-$i$ pinching, noncollapsing, and
canonical-neighborhood hypotheses before $T_i$, then it satisfies the strong
$(C,\epsilon)$-canonical-neighborhood assumption with parameter
$r_{i+1}$ throughout $[0,T)$.
\end{theorem}

\begin{proof}
\uses{def:strong-canonical-neighborhood-control,def:controlled-ricci-flow-with-surgery,thm:conditional-noncollapsing-extension,lem:canonical-neighborhood-first-failure,lem:rescaled-high-curvature-canonical-control,lem:first-failure-noncollapsing-bridge,lem:backward-parabolic-neighborhood-dichotomy,lem:uniform-backward-lifetime,thm:pointed-controlled-flow-compactness,thm:kappa-solution-canonical-neighborhoods,lem:bounded-flow-line-cap-exclusion}
If no pair $r_{i+1},\delta_{i+1}$ worked, choose
$r_a\downarrow0$ and, for each $a$, surgery controls
$\delta_{a,b}\downarrow0$ with counterexample flows.  Apply
\ref{lem:canonical-neighborhood-first-failure} and rescale each first-failure
point to scalar curvature one.  Pinching makes every smooth limit
nonnegatively curved.  The conditional theorem
\ref{thm:conditional-noncollapsing-extension} gives the stagewise volume-or-
positive-component alternative.  Apply
\ref{lem:first-failure-noncollapsing-bridge}: its canonical-neighborhood
outcome contradicts first failure, so its other outcome supplies one ordinary
noncollapsing constant on every fixed based ball.  The higher-curvature points
in the zero-time slices have doubled canonical control by
\ref{lem:rescaled-high-curvature-canonical-control}.

The canonical alternative in
\ref{lem:backward-parabolic-neighborhood-dichotomy} contradicts the choice
of basepoints, so compatible bounded-curvature backward neighborhoods exist
at every fixed radius.  By \ref{lem:uniform-backward-lifetime} they have a
common positive lifetime and can be extended maximally.  A finite maximal
lifetime would supply a bounded-curvature flow line into a surgery cap;
\ref{lem:bounded-flow-line-cap-exclusion} again contradicts first failure.
The lifetime is therefore infinite.

Apply \ref{thm:pointed-controlled-flow-compactness} and diagonalize over all
radii and backward times.  The limit is a $\kappa$-solution.
By \ref{thm:kappa-solution-canonical-neighborhoods}, its basepoint has a
strong $(C,\epsilon)$-canonical neighborhood, and the convergence clause of
that theorem transfers it to the first-failure points.  This final
contradiction proves that some $r_{i+1}$ and $\delta_{i+1}$ satisfy the
statement.
\end{proof}

\section{Surgery Continuation and Finiteness}

\begin{lemma}[Controlled restart after surgery]
\label{lem:controlled-surgery-restart}
\dcref{MT:ch15:15.8,MT:ch15:15.11}
\notready
Let a controlled Ricci flow reach a finite singular time $T$.  In its maximal
extension, cut every sufficiently curved horn at the central sphere of a
strong $\bar\delta(T)$-neck of scale $h(T)$, discard each horn end and each
component containing no controlled-curvature point, and insert the standard
caps.  For sufficiently small control, the post-surgery metric is smooth,
retains curvature pinching, satisfies the seven controlled-flow conditions,
and is initial data for a maximal ordinary Ricci flow.  The new flow retains
the stage noncollapsing and strong canonical-neighborhood estimates.
\end{lemma}

\begin{proof}
\uses{def:controlled-surgery-parameters,def:controlled-ricci-flow-with-surgery,def:maximal-ricci-flow-extension,thm:local-neck-surgery,thm:conditional-noncollapsing-extension,thm:strong-canonical-neighborhood-extension}
The maximal extension \ref{def:maximal-ricci-flow-extension} contains the
controlled-curvature region at time $T$.  The horn construction selects a
locally finite family of disjoint cutting necks, and compactness of the slice
makes the family finite.  Apply \ref{thm:local-neck-surgery} at each neck.
The metrics agree on the retained halves, the inserted caps are smooth, and
the local curvature comparison preserves pinching.  Discarding whole
components cannot spoil any estimate on the components which remain.

The exposed pieces are finitely many three-balls, their boundary spheres are
the prescribed central neck spheres, and the disappearing region is covered
by the canonical necks, caps, and discarded components.  Thus the seven
conditions in \ref{def:controlled-ricci-flow-with-surgery} persist.  Short-time
existence starts an ordinary flow from the compact post-surgery metric, and
uniqueness glues it to the continuing region.  Taking its maximal extension
gives the next smooth interval.  Finally
\ref{thm:conditional-noncollapsing-extension} and
\ref{thm:strong-canonical-neighborhood-extension} restore the two quantitative
stage estimates.
\end{proof}

\begin{theorem}[Map across a separating surgery]
\label{thm:separating-surgery-map}
\dcref{MT:ch15:15.12,MT:ch13:13.2}
\notready
Let $T$ be a surgery time of a controlled Ricci flow, and assume every
surgery sphere involved in a component $X(t)$ is separating.  For $t<T$
sufficiently close to $T$, let $X(T)$ be the component produced from
$X(t)$.  The points whose curvature remains bounded form an open set
$\Omega_t\subset X(t)$ on which the time-flow identifications converge to a
map onto the persistent open subset of $X(T)$.  For every $\eta>0$, after
choosing $t$ sufficiently close to $T$, the identification satisfies
\[
 (1+\eta)^{-1}d_{g(t)}(p,q)\le d_{g(T)}(\Phi_t p,\Phi_t q)
 \le(1+\eta)d_{g(t)}(p,q)
 \qquad(p,q\in\Omega_t).
\]
Assume that the radial maps on the removed sides agree with the persistent
identification on a collar and that the resulting collar interpolation is
$(1+\eta)$-Lipschitz.
It extends continuously to a map
\[
 F\colon X(t)\longrightarrow X(T)
\]
with $\operatorname{Lip}_{g(t),g(T)}(F)\le1+\eta$.  On each removed side,
the restriction is the standard radial collapse and hence is distance
nonincreasing there.
\end{theorem}

\begin{proof}
\uses{def:controlled-ricci-flow-with-surgery,thm:local-neck-surgery,lem:radial-surgery-map-nonexpanding}
On $\Omega_t$, curvature and all its derivatives remain bounded up to $T$;
the time flow therefore converges smoothly.  Uniform metric convergence on
compact subsets gives the displayed $(1+\eta)$ comparison, not an exact
isometry at the pre-surgery time.
Every component of the complement lies beyond one of the cutting necks.
Because each cutting sphere separates, it has a unique side not leading to
the selected persistent component.  Send that side through the radial
collapse to the corresponding new cap.  On the retained neck half this map
agrees with the persistent identification, so the local maps glue
continuously to $F$.

The map on each removed side is the radial map in
\ref{lem:radial-surgery-map-nonexpanding}; the metric on the persistent side
converges smoothly.  The explicit collar hypothesis supplies the only
remaining cross-neck estimate, and therefore the glued map has
$\operatorname{Lip}(F)\le1+\eta$.
\end{proof}

\begin{theorem}[Finite surgery count on bounded intervals]
\label{thm:finite-interval-surgery-bound}
\dcref{MT:ch17:17.12}
\notready
Let a controlled Ricci flow with surgery be defined on $[0,T)$ with
$T\le T_{i+1}$ and with the control functions fixed on that interval.  There
is
\[
 N=N\bigl(\operatorname{Vol}(M_0,g(0)),T_{i+1},r_{i+1},
             \bar\delta(T_{i+1})\bigr)<\infty
\]
such that the flow has at most $N$ surgery times in $[0,T)$.  In particular,
surgery times cannot accumulate in a bounded interval.
\end{theorem}

\begin{proof}
\uses{def:controlled-surgery-parameters,def:controlled-ricci-flow-with-surgery,lem:controlled-surgery-restart}
Away from surgery, curvature pinching gives $R\ge-6$ and hence
\[
 \frac d{dt}\operatorname{Vol}(M_t)
 =-\int_{M_t}R\,d\operatorname{vol}
 \le6\operatorname{Vol}(M_t).
\]
Thus the total volume immediately before every surgery is at most
\[
 e^{6T_{i+1}}\operatorname{Vol}(M_0)
\]
plus the earlier downward jumps.

A neck surgery removes a positive half-neck of scale $h(t)$ and inserts a
cap of volume at most $(1+\epsilon)Kh(t)^3$.  The choice of
$\bar\delta$ makes the removed cylindrical portion larger by at least
$c h(t)^3$.  Since $r$ and $\bar\delta$ are nonincreasing on the bounded
interval, the scale-selection function has a positive lower bound
\[
 h(t)\ge h\bigl(r_{i+1},\bar\delta(T_{i+1})\bigr)>0.
\]
Every neck surgery therefore removes at least one fixed volume $v_*>0$, so
the number of cutting operations is bounded by the exponential volume bound
divided by $v_*$.

If $S$ is the number of cutting operations and $D$ the number of whole
components discarded, the number of components is at most $N_0+S-D$.
Normalization gives a fixed positive lower volume for every initial
component, hence bounds $N_0$ by the initial volume.  It follows that $D$ is
bounded as well.  Each surgery time performs at least one cutting or discards
at least one component, which proves the stated bound.
\end{proof}

\section{All-Time Controlled Extension}

\begin{theorem}[All-time controlled Ricci flow with surgery]
\label{thm:all-time-controlled-surgery-flow}
\dcref{MT:ch15:15.9,MT:ch15:15.10}
\notready
Assume the conditional noncollapsing and strong canonical-neighborhood
extensions, together with their endpoint-transfer hypotheses, hold at every
finite stage.  Assume also that the local surgery construction supplies the
one-surgery connected-sum reconstruction contract at every restart.  There
are positive nonincreasing sequences
\[
 \mathbf K=(\kappa_i)_{i\ge0},\qquad
 \boldsymbol\delta=(\delta_i)_{i\ge0},\qquad
 \mathbf r=(r_i)_{i\ge0}
\]
such that every normalized compact three-manifold containing no embedded
$\mathbb{RP}^2$ with trivial normal bundle is the initial slice of a Ricci
flow with surgery on $[0,\infty)$.  For $t\in[T_i,T_{i+1})$, $x\in M_t$, and
$0<\rho\le\epsilon$, every parabolic ball satisfying the curvature hypothesis in
\ref{def:surgery-flow-noncollapsing} has volume at least
$\kappa_{i+1}\rho^3$ unless its center lies in a positively curved component.
The flow satisfies the strong $(C,\epsilon)$-canonical-neighborhood assumption
with parameter $r_{i+1}$, has curvature pinched toward positive, and uses
surgery control at most $\delta_{i+1}$.  Its surgery times are finite on every
bounded interval.
At every surgery time, the cutting and capping data also satisfy the
one-surgery connected-sum reconstruction contract, so the associated topology
change is part of the flow output.
\end{theorem}

\begin{proof}
\uses{def:surgery-flow-noncollapsing,def:controlled-surgery-parameters,def:controlled-ricci-flow-with-surgery,thm:normalized-short-time-ricci-flow,thm:conditional-noncollapsing-extension,thm:strong-canonical-neighborhood-extension,lem:controlled-surgery-restart,thm:finite-interval-surgery-bound}
Apply \ref{thm:normalized-short-time-ricci-flow} to the normalized initial
metric.  It gives the ordinary flow through $T_0=2^{-5}$, its initial
noncollapsing constant, and the curvature bounds which make the stage-zero
canonical and pinching requirements automatic after fixing the universal
constants.

Choose the sequences inductively.  Once the data through stage $i$ are fixed,
\ref{thm:strong-canonical-neighborhood-extension} supplies $r_{i+1}$ and an
upper surgery bound.  Decrease that bound if necessary so that
\ref{thm:conditional-noncollapsing-extension} supplies
$\kappa_{i+1}\le\kappa_i$.  This defines the next pieces of all three
control functions in \ref{def:controlled-surgery-parameters}.

Run the maximal ordinary flow from the current slice.  If it becomes singular
before $T_{i+1}$, apply \ref{lem:controlled-surgery-restart}; if no component
remains, declare every later slice empty.  Otherwise repeat from the new
compact slice.  The two extension theorems retain the quantitative hypotheses
after each restart.  By \ref{thm:finite-interval-surgery-bound}, only finitely
many restarts occur before $T_{i+1}$, so the construction reaches the stage
endpoint.

Induction over $i$ gives compatible controlled flows on all bounded time
intervals.  Their union is a Ricci flow with surgery on $[0,\infty)$ and has
every property asserted in the statement.
\end{proof}
